% Préparation d'une solution par dissolution d'un solide (L'essentiel, chapitre 1 ; exercice 6)
\begin{tikzpicture}[x=1cm,y=1cm,
  etape/.style={font=\otLight\fontsize{7.5}{9.2}\selectfont,text=otInk,align=center,text width=33mm,anchor=north}]
% 1. Pesée
\filldraw[pcAppareil,rounded corners=1mm] (-1.1,0) rectangle (1.1,0.55);
\filldraw[fill=otNavyDark,draw=otInk,line width=0.5pt,rounded corners=0.5mm] (-0.55,0.12) rectangle (0.55,0.42);
\node[text=white,font=\otMono\fontsize{8}{9}\selectfont] at (0,0.27) {2,50 g};
\draw[otInk!75,line width=0.6pt] (0,0.55) -- (0,0.72);
\filldraw[fill=otRule!40,draw=otInk!75,line width=0.6pt] (-0.85,0.72) rectangle (0.85,0.8);
\filldraw[fill=white,draw=otInk!75,line width=0.6pt] (-0.42,0.8) arc[start angle=180,end angle=360,x radius=0.42,y radius=0.16];
\foreach \x/\y in {-0.18/0.86,-0.05/0.9,0.08/0.87,0.2/0.85,0.0/0.95}{\fill[cyan!60!otNavy] (\x,\y) circle (0.035);}
\node[etape] at (0,-0.35) {\otRond{1}\enspace Peser la masse $m$ de solide dans une coupelle, sur une balance tarée.};
% 2. Introduction dans la fiole
\pic at (3.7,0) {fiole={hauteur=3.1,niveau=0}};
\draw[pcVerre] (3.36,2.95) -- (3.62,2.45) -- (3.62,2.15) (4.04,2.95) -- (3.78,2.45) -- (3.78,2.15);
\draw[pcVerre] (3.30,2.95) -- (4.10,2.95);
\foreach \x/\y in {3.66/0.1,3.72/0.12,3.78/0.1}{\fill[cyan!60!otNavy] (\x,\y) circle (0.035);}
\pic at (2.55,1.3) {pissette};
\node[etape] at (3.4,-0.35) {\otRond{2}\enspace L'introduire dans la fiole jaugée avec un entonnoir ; rincer coupelle et entonnoir à l'eau distillée.};
% 3. Dissolution
\pic at (7.0,0) {fiole={hauteur=3.1,niveau=0.55,liquide=cyan!25}};
\draw[otGreenDark,line width=0.7pt,-{Stealth[length=1.6mm]}] (7.95,1.0) arc[start angle=-40,end angle=70,radius=0.55];
\node[etape] at (7.0,-0.35) {\otRond{3}\enspace Remplir aux trois quarts d'eau distillée et agiter pour dissoudre tout le solide.};
% 4. Ajustement au trait de jauge
\pic at (10.4,0) {fiole={hauteur=3.1,niveau=1,liquide=cyan!25}};
\filldraw[pcAppareil] (10.27,3.1) rectangle (10.53,3.32);
\draw[pcRepere] (10.64,2.35) -- (11.25,2.75) node[pcLegende,anchor=west] {trait de jauge};
\node[etape] at (10.6,-0.35) {\otRond{4}\enspace Compléter jusqu'au trait de jauge, boucher, puis homogénéiser en retournant la fiole.};
\end{tikzpicture}
\par\smallskip\legende{Les quatre étapes d'une dissolution : la fiole jaugée fixe le volume de la solution avec précision. Schéma sans échelle.}
